% ===========================================================================
% Couverture — gabarit ESPRIT.
%
% Si l'école fournit un fichier de couverture officiel (.tex ou .docx converti),
% remplacez INTÉGRALEMENT cette page par le leur : la mise en page d'une
% couverture institutionnelle ne se reconstitue pas, elle se reprend.
% En attendant, cette page reprend l'ordonnancement attendu.
%
% Le logo se dépose dans assets/logos/esprit.png ; sans lui, la page compile
% quand même (un cadre discret prend sa place).
% ===========================================================================
\begin{titlepage}
\centering

\IfFileExists{assets/logos/esprit.png}{%
  \includegraphics[width=0.45\textwidth]{esprit}\par
}{%
  \fboxsep=10pt\fboxrule=0.4pt%
  \fcolorbox{piqbitExt}{white}{%
    \parbox{0.45\textwidth}{\centering\scriptsize\color{piqbitExt}
      Déposer le logo de l'établissement dans\\
      \texttt{assets/logos/esprit.png}}%
  }\par
}

\vspace{0.8cm}
{\large\bfseries \reportAcademicYear\par}
\vspace{0.3cm}
{\large PROJET DE FIN D'ÉTUDES\par}

\vspace{0.7cm}
% Le \parbox doit mesurer \textwidth MOINS la marge que \colorbox ajoute de
% chaque côté : \fboxsep vaut 10pt ici (fixé plus haut), donc un parbox à
% \textwidth produisait un bandeau 20pt trop large qui mordait sur la marge.
{\colorbox{piqbitInk}{%
  \parbox{\dimexpr\textwidth-2\fboxsep\relax}{\centering\vspace{6pt}
    {\Large\bfseries\color{white} \reportDegree}\vspace{6pt}}%
}\par}

\vspace{1.6cm}

{\large\bfseries Sujet :\par}
\vspace{0.4cm}
{\LARGE\bfseries \reportTitle\par}
\vspace{0.5cm}
{\large \reportSubtitle\par}

\vspace{1.8cm}

\begin{flushleft}
\hspace{2.5cm}\begin{minipage}{0.75\textwidth}
  \setlength{\parskip}{0.55em}
  {\large \textbf{Réalisé par :} \reportAuthor\par}
  {\large \textbf{Classe :} \reportClass\par}
  {\large \textbf{Entreprise :} \reportHostOrg\par}
  {\large \textbf{Encadrant professionnel :} \reportProfessionalSupervisor\par}
  {\large \textbf{Encadrant académique :} \reportAcademicSupervisor\par}
\end{minipage}
\end{flushleft}

\vfill
{\small \reportInstitution\par}
{\small Spécialité : \reportSpecialty\par}

\end{titlepage}
% `titlepage` REMET le compteur de page à 1 en se refermant (comportement de
% la classe standard). La couverture étant déjà la page i, la page suivante
% recevait elle aussi le numéro i : deux signets PDF portaient l'identifiant
% `page.i`, hyperref en ignorait un, et un lien pouvait mener à la mauvaise
% page. On reprend donc la numérotation là où elle doit être.
\setcounter{page}{2}
